\scanpage{51}\mspage{3\deleted{\unsure{9|8}}8}

et pour une flèche

\[
 \varnothing_{\unsure{M(I)}}\longrightarrow Y_\bullet
\]

Dire que cette flèche est dans $C_I$ signifie que les morphismes de transition $\deleted{\unsure{x}}Y_i\xrightarrow{u_i}Y_{i+1}$ sont dans $\deleted{\unsure{T}}C$, \uline{et} que les $Y_i$ sont $C$-cofibrants i.e. $\varnothing_M\longrightarrow Y_i$ est dans $C$, pour \unsure{tt} $i$. Mais la condition envisagée \unsure{en|dans} 2.1. \unsure{$1^\circ$|$2^\circ$}) est \deleted{\unsure{que}}\unsure{seulement} que les $u_i$ sont dans $C$. Donc en prenant pour $Y_\bullet$ un syst. inductif constant de valeur $Y$, avec $Y$ \unsure{non} $C$-cofibrant i.e. $\varnothing_M\longrightarrow Y$ non dans $C$, on a un contrexemple.

\unsure{Avec|Pour} le $=$ $I$, on voit que $u_\bullet:X_\bullet\longrightarrow Y_\bullet$ est dans $C$ ssi a) les $u_i\in C$, et b) les $u_{i,i+1}:Y_{i,i+1}\longrightarrow Y_{i+1}$ sont dans $C$. Si de plus les $u_i$ sont dans $\unsure{w|W}$, donc dans $TC$, on voit comme ci-dessus que $u_{i,i+1}\in TC_{\unsure{\varnothing}}$ pour tt $i$, donc $u_\bullet\in(TC)_I$. Donc ici on a encore

\[
 (TC)_I=C_I\cap W_{I^{\unsure{o|\circ}}},
\]

quoique $I$ ne soit pas artinien. J'ignore si l'hypothèse artinienne dans 2.1. $3^\circ$ est superflue ou non.

2) Soit $u_\bullet\in\operatorname{Fl}M(I^o)$, donc $u_\bullet$ est un objet de

\[
 \text{\uline{Fl}}\,M(I^o)\simeq M(I^o\times\unsure{\mathbb{\Delta}|\Delta}^{1})=M((I\times\mathbb{\Delta}^{1\,\mathrm{opp}})^{\mathrm{opp}}).
\]

\deleted{\unsure{P}} On vérifie que pour que cet objet soit $C$-cofibrant, il f. et s. que $1^\circ)$ $X_\bullet$ soit $C$-cofibrant, et $2^\circ)$ $u_\bullet\in C_I$ (ce qui implique que $Y_\bullet$ est aussi cofibrant)

% slanted
\margin{$I^o=I^{\mathrm{opp}}$}

% archive sheet number: 50

